\chapter{Travailler avec un assistant}

Un assistant fondé sur un modèle de langage lit des fichiers, lance des commandes, écrit du
texte et du code. Dans le projet qui sert d'exemple, il a transcrit des pages manuscrites,
écrit des preuves en Lean, cherché des références et rédigé des brouillons. Il s'agissait
de Claude, d'Anthropic, utilisé dans un terminal avec l'outil Claude Code. Ce chapitre décrit
l'organisation qui garde ce travail vérifiable.

\section{Déléguer, garder}

On délègue ce qui se vérifie : une compilation, une recherche dont on peut relire les
sources, une transcription qu'on peut comparer à la page, une preuve que le noyau de Lean
contrôle. On garde ce qui engage : le choix des questions, ce qu'on publie, ce qu'on écrit
sous son nom à d'autres personnes, et le jugement final sur la mathématique.

\begin{regle}
Tout ce que l'assistant produit est un brouillon jusqu'à ce qu'une vérification, par la
machine ou par une personne, l'ait établi. Le fichier dit lequel des deux.
\end{regle}

\section{Des procédures écrites}

Un assistant ne garde pas de mémoire d'une session à l'autre, sauf ce qu'on lui écrit. On
rédige donc les procédures récurrentes comme on rédigerait un protocole d'expérience : un
fichier par tâche, qui dit ce qu'il faut produire, dans quel ordre, avec quelles règles,
et ce qu'il ne faut jamais faire. Claude Code appelle ces fichiers des \emph{skills}.

\begin{exemple}
La procédure de transcription d'une page manuscrite tient en quelques pages. Elle fixe le
lot (vingt pages au plus par passe, parce qu'au-delà l'attention baisse sans prévenir), le
format (un sous-ensemble restreint de \LaTeX, avec des macros pour l'illisible et
l'incertain), l'en-tête (le modèle qui a lu, la date, la mention « non relu par une
personne »), et les interdits : corriger la mathématique de l'auteur, deviner un mot
illisible, transcrire une page sans mathématiques.
\end{exemple}

Ces procédures s'améliorent avec le travail. Quand une erreur revient, on ajoute une règle
et sa raison. Sans la raison, la règle tourne au rite, et personne n'ose plus la changer
quand elle ne sert plus.

\section{Une mémoire explicite}

Les décisions prises en route (ne pas citer telle source, ne jamais pousser de force sur
telle branche, contacter d'abord telle personne) se perdent si elles ne sont écrites nulle
part. On les garde dans des fichiers courts, un fait par fichier, avec la date et la raison,
que l'assistant relit au début de chaque session. Une décision qui change donne un nouveau
fichier, ou une correction datée de l'ancien.

\section{Déléguer en parallèle}

Un assistant peut confier des sous-agents à d'autres instances qui travaillent en même
temps : lire quatre documents à la fois, écrire deux preuves indépendantes. Le gain de temps
est réel ; le risque aussi. Chaque sous-agent reçoit une consigne complète, parce qu'elle ne
voit rien d'autre ; deux sous-agents ne touchent jamais les mêmes fichiers ; et leurs
résultats sont relus avant d'être intégrés.

\begin{warningblock}
Un rapport de sous-agent est écrit avec la même assurance, qu'il soit juste ou faux. On
vérifie les affirmations qui comptent (une référence, un énoncé, un résultat de
compilation) avant de les reprendre.
\end{warningblock}

\section{Des sessions dans le cloud}

Claude Code ne tourne pas seulement sur sa propre machine. On peut ouvrir une session dans
le cloud, depuis le site de Claude ou depuis l'application de bureau : une machine
distante clone le dépôt depuis GitHub, y travaille, et pousse ses commits. La machine
locale reste libre, et l'on peut fermer l'ordinateur pendant que la tâche avance.

C'est le bon outil pour une tâche longue et bien délimitée : une série de preuves à
compléter, une relecture de tout un dossier, la mise à jour des chiffres d'un article. Trois
précautions évitent la plupart des ennuis.

\begin{description}
  \item[La session ne voit que le dépôt.] Tout ce qui est ignoré par \fichier{.gitignore}
        (copies d'archives, transcriptions privées, brouillons locaux) lui manque. Une
        tâche qui en a besoin reste sur la machine locale, ou reçoit dans sa consigne ce
        qu'il lui faut, reformulé, jamais les fichiers privés eux-mêmes.
  \item[Une branche, un seul propriétaire.] Si une session dans le cloud et une session
        locale poussent sur la même branche, leurs commits se croisent et l'une écrase
        les hypothèses de l'autre. On attribue chaque branche à une seule session, et
        l'on récupère (\cmd{git fetch}) avant de reprendre la main.
  \item[La consigne tient seule.] Comme pour un sous-agent, la session distante ne
        connaît pas la conversation en cours. On lui écrit ce qu'il faut faire, où, avec
        quelles règles et comment vérifier, puis on relit ses commits comme ceux d'un
        collaborateur.
\end{description}

\begin{exemple}
Pendant qu'une session locale répondait aux reviews de plusieurs pull requests, une session dans le cloud a poussé sur les mêmes branches. Un de ses
commits cassait la compilation ; l'autre réintroduisait des citations que la review
venait de demander de retirer. La réparation a pris une heure. La règle « une branche, un
seul propriétaire » vient de là.
\end{exemple}

\section{Un budget}

Le travail d'un assistant consomme un quota de calcul. On fixe un plafond par semaine et
l'on vérifie où l'on en est avant de lancer une tâche longue. Le plafond oblige à choisir
ce qui se fait maintenant et ce qui attend.

\section{Une limite}

Je n'ai pas de bonne mesure de la fiabilité d'un travail produit ainsi. On peut compter les
points relus par une personne et ceux qu'elle a corrigés, et en tirer une estimation sur
ces échantillons ; rien de plus. Un article honnête sur un tel projet le dit, et ne
généralise pas.

\section{Déclarer}

Les revues et les sites d'entraide ont des politiques différentes sur les textes écrits avec
un modèle de langage. Certaines revues les refusent ; certains sites interdisent les
réponses générées ; d'autres demandent seulement une déclaration. On lit la politique avant
de soumettre, on déclare l'usage dans les termes demandés, et l'on écrit soi-même ce qu'on
publie sous son nom.
